\section{电磁感应}

\subsection{电磁感应的基本规律}

\subsubsection{电源和电动势}

电源通过非静电力对电荷产生作用力。定义非经典力场强为电源对单位正电荷产生的作用力，即：
$$
\vect{K} = \frac{\vect{F}_{\text{非}}}{q}
$$

定义电动势为电源对单位正电荷所作的功，即：
$$
\varepsilon = \frac{A_{\text{非}}}{q} = \int_{-}^{+} \vect{K} \vdot \dd{\vect{l}}
$$

\subsubsection{电磁感应现象}

变化的磁场可以在回路中产生感应电动势。

\subsubsection{法拉弟定律}

\textbf{法拉第定律}指出，闭合回路中的电磁感应电动势正比于回路磁通量的变化，即：
$$
E = \dv{\Phi}{t} = \dv{t} \iint_C \vect{B} \vdot \dd{\vect{S}}
$$

\subsubsection{楞次定律}

\textbf{楞次定律}表述为：闭合回路中感应电流的方向，总是使它所激发的磁场来阻止引起感应电流的磁场变化。

\subsection{作业}

\begin{problem}
	习题 4.18
	\begin{solution}
		线圈的磁矩为：
		$$
		\abs{\vect{p}_m} = I S = \frac{\pi}{20} \,\mathrm{A \cdot m^2}
		$$
		因此所受磁力矩为：
		$$
		\abs{\vect{M}} = \abs{\vect{p}_m} \cdot \abs{\vect{B}} = \frac{\pi}{40} \,\mathrm{N \cdot m}
		$$
		根据叉积判断，磁力矩方向为竖直向上。
	\end{solution}
\end{problem}

\begin{problem}
	习题 4.20
	\begin{solution}
		\begin{enumerate}
			\item[\textbf{1)}] 使用微元分析，对于一个圆心角为 $\dd{\theta}$，弧长为 $\dd{l} = d \dd{\theta}$ 的微元，其收到两侧的张力，张力和安培力平衡，有
			$$
			2 F \sin \frac{\dd{\theta}}{2} = B I \dd{l}
			$$
			解得 $F = B I d$。因此每单位截面积上张力为 $p = \frac{F}{S} = 1 \times 10^{6} \,\mathrm{Pa}$。

			\item[\textbf{2)}] $p_0 = 1.96 \times 10^{6} \,\mathrm{Pa} > p$，因此不会断裂。
		\end{enumerate}
	\end{solution}
\end{problem}

\begin{problem}
	习题 4.22
	\begin{solution}
		\begin{enumerate}
			\item[\textbf{1)}]
			$$
			\begin{dcases}
				\frac{1}{2} m v^2 = q U \\
				q v B = \frac{m v^2}{R} \\
				R = \frac{x}{2}
			\end{dcases}
			$$
			解得 $m = \frac{8 q B^2 x^2}{U}$。

			\item[\textbf{2)}] 从 \textbf{1)} 得到 $\frac{q}{m} = \frac{U}{8 B^2 x^2}$。代入数据得到 $\frac{q}{m} \approx 6.88 \times 10^{4} \,\mathrm{C / kg}$。

			\item[\textbf{3)}] 从 \textbf{1)} 得到 $x = \sqrt{\frac{m U}{8 q B^2}}$。代入数据得到 $x \approx 0.03 \,\mathrm{m}$。
		\end{enumerate}
	\end{solution}
\end{problem}

\begin{problem}
	习题 4.23
	\begin{solution}
		\begin{enumerate}
			\item[\textbf{1)}] 平衡时，对于一个电子，其收到的洛伦兹力和电场力平衡：
			$$
			e v B = \frac{e U_{AA'}}{b}
			$$
			而根据电流的定义，有
			$$
			I = n e S v \implies v = \frac{I}{n e b d}
			$$
			代入得到 $U_{AA'} = \frac{I B}{n e d} \approx 2.2 \times 10^{-5} \,\mathrm{V}$

			\item[\textbf{2)}] 根据式子，是无关的。
		\end{enumerate}
	\end{solution}
\end{problem}

\begin{problem}
	习题 5.3
	\begin{solution}
		\begin{enumerate}
			\item[\textbf{1)}] 根据环路定理，磁场强度分布：
			$$
			H(r) = \begin{dcases}
				\frac{I r}{2 \pi R_1^2} &,\ 0 \le r < R_1 \\
				\frac{I}{2 \pi r} &,\ r \ge R_1 \\
			\end{dcases}
			$$
			根据 $H$ 和 $B$ 的关系，得到：
			$$
			B(r) = \begin{dcases}
				\frac{\mu_0 I r}{2 \pi R_1^2} &,\ 0 \le r < R_1 \\
				\frac{\mu_0 \mu_r I}{2 \pi r} &,\ R_1 \le r < R_2 \\
				\frac{\mu_0 I}{2 \pi r} &,\ r \ge R_2
			\end{dcases}
			$$
			其图像如下：

			\begin{figure}[H]
				\centering
				\begin{minipage}{0.48\textwidth}
					\centering
					\begin{tikzpicture}[>=latex,x=1cm,y=1cm,every node/.style={font=\scriptsize}]
						\def\Rone{2.0}
						\def\xmax{5.4}
						\def\ymax{3.2}
						\def\Hone{2.4}

						\draw[->] (0,0) -- (\xmax,0) node[below right] {$r$};
						\draw[->] (0,0) -- (0,\ymax) node[above left] {$H(r)$};

						\draw (\Rone,0) -- ++(0,-0.08) node[below] {$R_1$};
						\draw (0,\Hone) -- ++(-0.08,0) node[left] {$\dfrac{I}{2\pi R_1}$};

						\draw[dashed] (\Rone,0) -- (\Rone,\Hone);
						\draw[dashed] (0,\Hone) -- (\Rone,\Hone);

						\draw[thick] (0,0) -- (\Rone,\Hone);
						\draw[thick,domain=\Rone:\xmax-0.2,samples=100,smooth]
							plot (\x, {\Hone*(\Rone/\x)});

						\fill (\Rone,\Hone) circle (2pt);
						\node at (3.1,2.85) {$H(r)$};
					\end{tikzpicture}
				\end{minipage}
				\hfill
				\begin{minipage}{0.48\textwidth}
					\centering
					\begin{tikzpicture}[>=latex,x=1cm,y=1cm,every node/.style={font=\scriptsize}]
						\def\Rone{1.6}
						\def\Rtwo{4.0}
						\def\xmax{5.8}
						\def\ymax{5.4}
						\def\Bone{1.2}
						\def\Bmuone{4.8}
						\def\Bmutwo{1.92}
						\def\Btwo{0.48}

						\draw[->] (0,0) -- (\xmax,0) node[below right] {$r$};
						\draw[->] (0,0) -- (0,\ymax) node[above left] {$B(r)$};

						\draw (\Rone,0) -- ++(0,-0.08) node[below] {$R_1$};
						\draw (\Rtwo,0) -- ++(0,-0.08) node[below] {$R_2$};

						\draw (0,\Bone) -- ++(-0.08,0) node[left] {$\dfrac{\mu_0 I}{2\pi R_1}$};
						\draw (0,\Bmuone) -- ++(-0.08,0) node[left] {$\dfrac{\mu_0\mu_r I}{2\pi R_1}$};
						\draw (0,\Bmutwo) -- ++(-0.08,0) node[left] {$\dfrac{\mu_0\mu_r I}{2\pi R_2}$};
						\draw (0,\Btwo) -- ++(-0.08,0) node[left] {$\dfrac{\mu_0 I}{2\pi R_2}$};

						\draw[dashed] (\Rone,0) -- (\Rone,\Bmuone);
						\draw[dashed] (\Rtwo,0) -- (\Rtwo,\Bmutwo);

						\draw[thick] (0,0) -- (\Rone,\Bone);
						\draw[thick,domain=\Rone:\Rtwo,samples=100,smooth]
							plot (\x, {\Bmuone*(\Rone/\x)});
						\draw[thick,domain=\Rtwo:\xmax-0.2,samples=100,smooth]
							plot (\x, {\Btwo*(\Rtwo/\x)});

						\draw[fill=white,thick] (\Rone,\Bone) circle (2pt);
						\fill (\Rone,\Bmuone) circle (2pt);
						\draw[fill=white,thick] (\Rtwo,\Bmutwo) circle (2pt);
						\fill (\Rtwo,\Btwo) circle (2pt);

						\node at (3.0,5.05) {$B(r)$};
					\end{tikzpicture}
				\end{minipage}
				\caption{$H(r)$ 和 $B(r)$ 的图像}
			\end{figure}
		\end{enumerate}
	\end{solution}
\end{problem}

\begin{problem}
	习题 5.5
	\begin{solution}
		\begin{enumerate}
			\item[\textbf{1)}] $\abs{\vect{B}} = \frac{\Phi}{S} = 2.0 \times 10^{-2} \,\mathrm{T}$；
			\item[\textbf{2)}] 选择铁环中心线应用环路定理：
			$$
			\abs{\vect{H}} \cdot l = N I \implies \abs{\vect{H}} = \frac{N I}{l} = 32 \,\mathrm{A / m}
			$$
			\item[\textbf{3)}] 根据磁感应强度和磁场强度的关系：
			$$
			\vect{B} = \mu_0 (\vect{H} + \vect{M}) \implies \abs{\vect{M}} = \frac{\abs{\vect{B}}}{\mu_0} - \abs{\vect{H}} \approx 1.59 \times 10^{-4} \,\mathrm{A / m}
			$$
			\item[\textbf{4)}]
			$$
			\chi_m = \frac{\abs{\vect{M}}}{\abs{\vect{H}}} \approx 4.96 \times 10^{2},\quad \mu_r = 1 + \chi_m \approx 4.97 \times 10^{2}
			$$
		\end{enumerate}
	\end{solution}
\end{problem}

\begin{problem}
	习题 5.6
	\begin{solution}
		\begin{enumerate}
			\item[\textbf{1)}] 取铁环中线应用环路定理：
			$$
			\abs{\vect{H}} \cdot l = N I l \implies \abs{\vect{H}} = N I = 2.0 \times 10^{3} \,\mathrm{A / m}
			$$
			\item[\textbf{2)}]
			$$
			\abs{\vect{M}} = \abs{\frac{\vect{B}}{\mu_0} - \vect{H}} \approx 7.94 \times 10^{5} \,\mathrm{A / m}
			$$
			\item[\textbf{3)}]
			$$
			\mu_r = \frac{\abs{\vect{B}}}{\mu_0 \abs{H}} \approx 4.0 \times 10^{2}
			$$
		\end{enumerate}
	\end{solution}
\end{problem}

\begin{problem}
	习题 6.1
	\begin{solution}
		$$
		E = \dv{\Phi}{t} = \dv{\ab(B (d_0 + v t) l \cos \alpha)}{t} = B v l \cos \alpha \approx 1.5 \,\mathrm{V}
		$$
		电流方向为逆时针。
	\end{solution}
\end{problem}

\begin{problem}
	习题 6.3
	\begin{solution}
		\begin{enumerate}
			\item[\textbf{1)}]
			$$
			\Phi = \int_{a}^{b} \frac{\mu_0 i}{2 \pi x} l \dd{x} = \frac{I_0 l \sin \omega t}{2 \pi} \ln \frac{b}{a}
			$$
			\item[\textbf{2)}]
			$$
			E = \dv{\Phi}{t} = \frac{\omega I_0 l \cos \omega t}{2 \pi} \ln \frac{b}{a}
			$$
		\end{enumerate}
	\end{solution}
\end{problem}

\begin{problem}
	习题 6.4
	\begin{solution}
		\begin{enumerate}
			\item[\textbf{1)}] 在中轴线上的磁感应强度为：
			$$
			\begin{aligned}
				\vect{B} & = \oint_L \frac{\mu_0}{4 \pi} \frac{I \dd{\vect{l}} \cp \vect{r}}{r^3} \\
				& = \vect{e}_z \oint_L \frac{\mu_0}{4\pi} \frac{I \dd{l}}{R^2 + x^2} \cdot \frac{1}{\sqrt{R^2 + x^2}} \\
				& = \frac{\mu_0 I R^2}{2 \ab(R^2 + x^2)^{\frac{3}{2}}} \vect{e}_z
			\end{aligned}
			$$
			因此磁通量为 $\Phi = \pi r^2 B = \frac{\pi \mu_0 I R^2 r^2}{2 \ab(R^2 + x^2)^{\frac{3}{2}}}$。

			\item[\textbf{2)}]
			$$
			\dv{x}{t} = v \implies E = \dv{\Phi}{t} = -\frac{3 \pi \mu_0 I R^2 r^2 v x}{2 \ab(R^2 + x^2)^{\frac{5}{2}}}
			$$
			方向顺时针。
		\end{enumerate}
	\end{solution}
\end{problem}

\begin{problem}
	习题 6.5
	\begin{solution}
		\begin{itemize}
			\item 当 $0 \le r < R$ 时：
			$$
			E = \frac{U}{2 \pi r} = \dv{B}{t} \frac{\pi r^2}{2 \pi r} = \frac{r}{2} \dv{B}{t}
			$$
			\item 当 $r \ge R$ 时：
			$$
			E = \frac{U}{2 \pi r} = \dv{B}{t} \frac{\pi R^2}{2 \pi r} = \frac{R^2}{2 r} \dv{B}{t}
			$$
		\end{itemize}
		代入数据得到答案为 $5 \times 10^{-4} \,\mathrm{V / m},\ 1.25 \times 10^{-3}\,\mathrm{V / m},\ 2.5 \times 10^{-3} \,\mathrm{V / m},\ 1.25 \times 10^{-3}\,\mathrm{V / m}$。
	\end{solution}
\end{problem}